\documentclass[10pt,a4paper]{amsart}
\usepackage[margin=19mm]{geometry}
\usepackage{amsmath,amssymb,mathtools}
\usepackage[scr=rsfs,cal=euler]{mathalfa}
\usepackage{xfrac}
\usepackage[unicode,hidelinks,bookmarks=false]{hyperref}
\newcommand{\ie}{i.e.\ }
\newcommand{\cf}{cf.\ }
\newcommand{\resp}{respectively\ }
\newcommand{\Subsection}{\subsection}
\providecommand{\nobreakdash}{\nobreak\hbox{-}}
\setlength{\emergencystretch}{2em}
\multlinegap0pt
\usepackage{bm}
\usepackage{bbm}
\usepackage{dsfont}
\usepackage{xspace}
\usepackage{cancel}
\usepackage{mathtools}
\usepackage{amsthm}
\makeatletter
\@ifundefined{Thm}{\newtheorem{Thm}{Thm}}{}
\@ifundefined{Prop}{\newtheorem{Prop}[Thm]{Proposition}}{}
\makeatother
\title[Supercritical sharpness for the random cluster representatio]{Supercritical sharpness for the random cluster representation of real-valued spin models}
\author{Trishen S. Gunaratnam, Dmitry Krachun, Christoforos Panagiotis, Romain Panis, Franco Severo}
\date{Source: arXiv:2608.18045v1 (CC BY 4.0)}
\begin{document}
\maketitle

\noindent\emph{Source and licence.} Supercritical sharpness for the random cluster representation of real-valued spin models. Version arXiv:2608.18045v1 (CC BY 4.0), distributed under the Creative Commons Attribution 4.0 International licence, as recorded in the arXiv licence metadata for this exact version and shown on the abstract page https://arxiv.org/abs/2608.18045v1. Full rights notice: licenses/arxiv-2608.18045-CC-BY-4.0.txt.\par\smallskip

\noindent\textit{Editorial scope.} Counted unit: the Prop whose source label is \texttt{prop: lebowitz} in the article (source0.tex lines 2059--2065, source label \texttt{prop: lebowitz}), delivered together with the article's own complete original proof of it (source0.tex lines 2143--2222, one \texttt{proof} environment of 80 lines). No mathematics is written, repaired or re-derived: statement and proof are the article's own text line for line. The proof is detached from the statement in the source: the article states the result at lines 2059--2065 and prints its proof at lines 2143--2222, and the proof environment's own header names the statement, so the association is the article's own. Required context retained uncounted: none. Every definition, display and label of those lines is retained unchanged. Omitted from this package and not redistributed: 1--2058, 2066--2142, 2223--2513. Nothing inside a retained excerpt is omitted or altered: every retained body is delivered exactly as the article's lines are written, so no omission receipt of any kind accompanies this package. Citations: the retained excerpts carry no citation command at all, so no bibliography entry of the article is transcribed and no reference is redistributed. Reference adaptation: none was needed and none was made; every reference inside the delivered unit resolves inside it, so no unresolved reference or citation is produced. Printed slips kept literal: no printed word, symbol, hypothesis or number of the counted statement or of its proof was altered. Layout and class: the article's own class is \texttt{article} with packages including enumitem, geometry, footmisc, amsfonts, amsthm, amsmath, amssymb, centernot, fontenc, cite, mathrsfs, amscd, inputenc, t1enc; the wrapper is the lane's pinned \texttt{amsart} editorial wrapper at 10pt on A4, which restates the theorem-like environments the retained text needs and adds the article's own macro definitions that the retained text needs (none), with extra wrapper packages (bm, bbm, dsfont, xspace, cancel, mathtools). The delivered numbering is the unit's own. Assets: no retained span contains a figure environment, a graphics-inclusion command, a PDF-page inclusion, a table or a TikZ picture; no image file is copied into this package, the package declares no asset dependency and no image file is redistributed with it. Licence: version arXiv:2608.18045v1 of this article is distributed under the Creative Commons Attribution 4.0 International licence, as recorded in the arXiv licence metadata for this exact version and shown on the abstract page https://arxiv.org/abs/2608.18045v1. A full rights notice is delivered at licenses/arxiv-2608.18045-CC-BY-4.0.txt. Characters: every retained excerpt is pure ASCII, so the delivered engine, XeLaTeX with its default Latin Modern text font, reports no missing character.
\begin{Prop} \label{prop: lebowitz}
Fix a finite graph $G = (V,E)$. Consider two sets of coupling constants $\{J_{xy}\}_{xy\in E}$ and $\{J'_{xy}\}_{xy\in E}$ satisfying $|J'_{xy}| \le J_{xy}$ for every $xy\in E$, and two external magnetic fields $\mathsf{h}$ and $\mathsf{h}'$ satisfying $|\mathsf{h}'_x|\le \mathsf{h}_x$ for every $x\in V$. Consider a family of single-site measures $(\rho_x)_{x\in V}$ such that for every $x\in V$, $\rho_x$ is either an even measure over $\mathbb{R}$ or $\rho_x$ is supported on $\mathbb{R}^+$. For every pair of moment functions $A, B:V\to\mathbb{N}$, one has
\begin{equation}
\langle \tau_A \tau_B \rangle - \langle \tau_A \tau_B \rangle' \ge |\langle \tau_A \rangle \langle\tau_B \rangle' - \langle \tau_A \rangle' \langle \tau_B \rangle|,
\end{equation}
where $\langle \cdot \rangle$ and $\langle \cdot \rangle'$ stand for the expectations under the measures $\nu_{G,\rho,J, \mathsf{h}}$ and $\nu_{G,\rho,J',\mathsf{h}'}$ respectively\footnote{Note that $\beta$ has been chosen equal to $1$ here which explains why it disappeared from the notation.}.
\end{Prop}

\begin{proof}[Proof of Proposition~\textup{\ref{prop: lebowitz}}]
Let $\tau, \tau'$ be configurations sampled independently according to the measures $\mu:=\nu_{G,\rho,J,\mathsf{h}}$ and $\mu':=\nu_{G,\rho,J',\mathsf{h}'}$.
It suffices to prove that
\begin{equation}\label{eq:pfleb1}
\mathbb{E}_{\mu\otimes\mu'} [\tau_A\tau_B-\tau_A'\tau_B'] \ge \mathbb{E}_{\mu\otimes\mu'}[\tau_A\tau_B'-\tau_A'\tau_B],
\end{equation}
because the right-hand side equals $\langle \tau_A \rangle \langle\tau_B \rangle' - \langle \tau_A \rangle' \langle \tau_B \rangle$, and the second required inequality follows from exchanging the roles of $A$ and $B$.
Observe that \eqref{eq:pfleb1} can be rewritten as
\begin{equation}\label{eq:pfleb2}
\mathbb{E}_{\mu\otimes\mu'} [(\tau_A + \tau'_A)(\tau_B - \tau'_B)] \ge 0.
\end{equation}
To show \eqref{eq:pfleb2}, we introduce new spin variables labelled by the set of vertices $V$:
\begin{equation}
u_x := \tau_x + \tau'_x, \qquad v_x := \tau_x - \tau'_x.
\end{equation}
In terms of these new variables, we can write 
\begin{equation}
\tau_A+\tau_A' = \prod_{x\in V} \left(\frac{u_x+v_x}{2}\right)^{A_x} + \prod_{x\in V} \left(\frac{u_x-v_x}{2}\right)^{A_x}=2^{-|A|}\sum_{S\le A}\binom{A}{S}(1+(-1)^{|A|-|S|})u_S v_{A\setminus S},
\end{equation}
and
\begin{equation}
\tau_B-\tau_B' = \prod_{y\in V} \left(\frac{u_y+v_y}{2}\right)^{B_y} - \prod_{y\in V} \left(\frac{u_y-v_y}{2}\right)^{B_y}=
2^{-|B|}\sum_{S\le B}\binom{B}{S}(1-(-1)^{|B|-|S|})u_S v_{B\setminus S},
\end{equation}
where $|A|:=\sum_{x\in V} A_x$, the notation $S\le A$ means that $S:V \to\mathbb{N}$ is a moment function such that $S_x\le A_x$ for every $x\in V$, $A\setminus S$ denotes the moment function which is equal to $A_x- S_x$ for every $x\in V$, and $\binom{A}{S}:=\prod_{x\in V}\binom{A_x}{S_x}$.
Hence, it suffices to show that for all moment functions $A,B$,  
\begin{equation}\label{eq:what-we-need-to-prove}
\mathbb{E}_{\mu\otimes\mu'}[u_Av_B]\geq 0.
\end{equation}
We now rewrite the measure $\mu\otimes \mu'$ in terms of the new spin variables $u,v$, the new coupling constants 
\begin{equation}
K_{xy} := \frac{1}{2}(J_{xy} + J'_{xy}), \qquad L_{xy} := \frac{1}{2}(J_{xy} - J'_{xy}),
\end{equation}
(both of which are non-negative by the assumption on $J_{xy}, J'_{xy}$), as well as the new external magnetic fields
\begin{equation}
K_x:=\frac{1}{2}(\mathsf{h}_x+\mathsf{h}'_x), \qquad L_x:= \frac{1}{2}(\mathsf{h}_x-\mathsf{h}'_x)
\end{equation}
(both of which are non-negative by the assumption on $\mathsf{h}_x,\mathsf{h}'_x$).
We have 
\begin{align*}
   \mathrm{d}\mu\otimes \mu'[(\tau, \tau')] \propto \exp\left\{ \sum_{xy\in E} (J_{xy} \tau_x \tau_y +J'_{xy}\tau_x'\tau_y') +\sum_{x\in V} (\mathsf{h}_x\tau_x+\mathsf{h}_x'\tau_x') \right\}\prod_{x\in V}\mathrm{d}\rho_x(\tau_x)\mathrm{d}\rho_x(\tau'_x),
\end{align*}
and using the new variables we can write 
\begin{align*}
J_{xy} \tau_x \tau_y +J'_{xy}\tau_x'\tau_y'
&=
(K_{xy}+L_{xy})\cdot \frac{u_x+v_x}{2}\cdot\frac{u_y+v_y}{2}+(K_{xy}-L_{xy})\cdot \frac{u_x-v_x}{2}\cdot\frac{u_y-v_y}{2}
\\&=
K_{xy}\cdot \frac{u_xu_y+v_xv_y}{2}+L_{xy}\cdot \frac{u_xv_y+v_xu_y}{2}
\end{align*}
and
\begin{align*}
\mathsf{h}_x \tau_x +\mathsf{h}'_x\tau_x'
&=
(K_x+L_x)\cdot \frac{u_x+v_x}{2}+(K_x-L_x)\cdot \frac{u_x-v_x}{2}=
K_x\cdot u_x+L_x\cdot v_x.
\end{align*}
Also, let $\nu_x$ denote the pushforward of $\rho_x\times \rho_x$ by the map $F(a,b)=\left(a+b,a-b\right)$.
Since $F$ is a continuous bijection, the whole measure can now be written as 
\begin{align}
\frac{\mathrm{d}\mu\otimes \mu'[(u, v)]}{\prod_{x\in V}\mathrm{d}\nu_x(u_x,v_x)}\propto \exp\Biggl\{\sum_{xy\in E} \left(K_{xy} \frac{u_xu_y+v_xv_y}{2}+L_{xy} \frac{u_xv_y+v_xu_y}{2}\right) +\sum_{x\in V}(K_xu_x+L_xv_x)\Biggr\}.
\end{align}



We say that a measure $\nu$ on $\mathbb{R}^2$ is even if for every Borel measurable set $A$, we have $\nu(f_1(A))=\nu(f_2(A))=\nu(A)$, where $f_1(x,y)=(-x,y)$ and $f_2(x,y)=(x,-y)$. We claim that for every $x\in V$, the above measure $\nu_x$ is either an even measure on $\mathbb{R}^2$ or it is supported on $\mathbb{R}^+\times \mathbb{R}$ and for every Borel measurable set $A$, $\nu_x(f_2(A))=\nu_x(A)$. Indeed, it suffices to consider sets of the form $A_1\times A_2$, since these sets generate the $\sigma$-algebra. If $\rho_x$ is an even measure, then $\rho_x\otimes \rho_x$ is an even measure on $\mathbb{R}^2$ and invariant under $f_3(x,y)=(-y,-x)$. Hence
\begin{equation}\nu_x(A_1,A_2)=\rho_x\otimes \rho_x(F^{-1}(A_1,A_2))=\rho_x\otimes \rho_x(f_3(F^{-1}(A_1,A_2)))=\nu_x(-A_1,A_2),
\end{equation}
where we used that $F^{-1}(a,b)=(\tfrac{a+b}{2},\tfrac{a-b}{2})$ and $f_3(F^{-1}(a,b))=(\tfrac{-a+b}{2},\tfrac{-a-b}{2})$.
If $\rho_x$ is supported on $\mathbb{R}^+$, then $\nu_x$ is supported on $\mathbb{R}^+\times \mathbb{R}$ since $F^{-1}(A_1,A_2)\subset (\mathbb{R}^+)^2$ implies that $A_1\subset \mathbb{R}^+$.
On the other hand, in any case, by the invariance of $\rho_x\otimes\rho_x$ under $f_4(x,y)=(y,x)$,
\begin{equation}
\nu_x(A_1,A_2)=\rho_x\otimes\rho_x(F^{-1}(A_1,A_2))=\rho_x\otimes\rho_x(f_4(F^{-1}(A_1,A_2)))=\nu_x(A_1,-A_2).
\end{equation}
This proves the claim.

We can view the latter expression as a spin measure on the enhanced graph $G\times K_2$ with vertex set $V^2:=\{(x,i): x\in V, i\in \{1,2\}\}$ and edge set $E^2=\{\{(x,i),(y,j)\}: xy\in E, i,j\in \{1,2\}\}$. Let us write for $x\in V$, $\tau_{(x,1)}=u_x$ and $\tau_{(x,2)}=v_x$. Then, by the above property of each $\nu_x$, conditioning on the absolute value field $(|\tau_z|)_{z\in V^2}$, the sign field $(\textup{sgn}(\tau_z))_{z\in V^2}$ is distributed according to a (ferromagnetic) Ising measure on $G\times K_2$ with coupling constants $\tfrac{1}{2} K_{xy}\cdot |\tau_{(x,i)}|\cdot |\tau_{(y,j)}|$ for $i=j$ and $\tfrac{1}{2}L_{xy}\cdot |\tau_{(x,i)}|\cdot |\tau_{(y,j)}|$ for $i\neq j$, external magnetic field $K_x|\tau_{(x,1)}|$ and $L_x|\tau_{(x,2)}|$, and non-negative boundary conditions. Observe that when $x\in V$ is such that $\rho_x$ is supported on $\mathbb R^+$, one has $\textup{sgn}(\tau_{(x,1)})=1$.

Hence, the inequality \eqref{eq:what-we-need-to-prove} follows from Griffiths' first inequality for this Ising measure. This completes the proof.
\end{proof}

\end{document}
